\section{서론}
\label{sec:서론}

무인 비행체 애드혹 네트워크(FANET)는 재난 지역의 긴급 통신망 구축, 환경 감시 시스템, 신속하게 배치 가능한 전술 군사 네트워크 등 고정 지상 통신 인프라가 부재한 상황에서 무선 멀티홉 네트워크를 형성하는 핵심 기술로 주목받고 있다. FANET에서 개별 무인기(UAV)는 패킷을 목적지까지 전달하는 라우터 역할을 유기적으로 수행한다. 그러나 dynamic한 3차원 UAV 이동 환경과 이로 인한 빈번한 링크 단절, 채널 페이딩, 신호 감쇠 현상은 신뢰성 있는 라우팅을 설계하는 데 본질적인 장벽이 된다. 더욱이 UAV는 비행 시간, 가벼운 하드웨어 중량 등 내재적인 컴퓨팅 및 전력 자원의 한계를 지니고 있다. 결과적으로 실용적인 FANET 라우팅은 다음 두 가지 상충하는 요구사항을 만족해야 한다. 첫째, 토폴로지와 링크의 불안정성을 조기에 판단하여 끊김 없는 신뢰성 있는 경로를 예측해야 하며, 둘째, 통신 및 계산 오버헤드가 극도로 적어 개별 UAV에서 완벽히 분산 실행이 가능해야 한다.

기존 모바일 애드혹 네트워크(MANET) 및 차량 애드혹 네트워크(VANET) 환경에서 제안된 AODV나 OLSR과 같은 토폴로지 기반 라우팅 프로토콜은 UAV의 3차원 고속 이동 환경에서 경로 탐색 및 유지보수를 위한 제어 메시지 오버헤드가 비대해져 성능이 급격히 저하된다. 이러한 대안으로 제시된 지리적 라우팅(Geographic Routing) 기법인 GPSR~\cite{gpsr}은 전역 라우팅 테이블을 유지하지 않고 오직 목적지와 1-hop 이웃 노드들의 위치 정보만을 바탕으로 greedy하게 차홉 노드를 선정한다. 이 기법은 통신 오버헤드가 적어 대규모 군집에 적합하지만, 목적지 방향에 더 이상 가까운 이웃이 존재하지 않아 패킷이 갇히는 라우팅 홀(Routing Hole) 문제에 매우 취약하다. 또한 greedy한 지리적 진행만을 고려하므로 링크 수명, 버퍼 혼잡도 등을 무시해 토폴로지 변화 시 잦은 패킷 유실을 야기한다.

이를 극복하기 위해 최근 다중 에이전트 강화학습(MARL) 기반의 라우팅 기법들이 활발히 연구되고 있다. 강화학습 기법은 패킷 전송 성공률, 지연 시간, 대기열(Queue) 상태 등을 보상 함수로 정식화하여 동적으로 경로를 최적화한다. 그러나 고성능의 MARL 모델은 전역 상태 정보를 교환해야 하거나, 분산 실행 시에도 지속적인 에이전트 간 제어 메시지 전송(DRAMA~\cite{drama} 등)을 동반한다. 이는 대역폭이 좁고 잡음이 심한 실제 FANET 배포 환경에서 실행 불가능한 '배포 격차(Deployment Gap)'를 유발한다.

본 논문에서는 이러한 한계를 메우기 위해 글로벌-로컬 정책 증류(Policy Distillation) 및 예측적 위험 스위칭(Predictive Risk Switching) 기반의 저지연·고강인성 라우팅 프레임워크인 \method를 제안한다. 제안 기법은 중앙 집중식 학습 및 분산 실행(CTDE) 패러다임 하에서 오프라인 학습 단계 시 전역 그래프 정보에 접근 가능한 특권적 교사(Privileged Teacher) 모델을 PPO 알고리즘으로 학습시킨다. 이후 교사의 행동 확률 분포를 1-hop 로컬 관측 정보만을 입력으로 받는 경량 학생(Student) MLP 정책에 지식 증류하여 지식 이전을 달성한다. 실제 배포 단계에서는 예측적 위험 스위칭 기법을 적용한다. 평상시(Benign case)에는 가벼운 nominal 지리적 라우팅을 수행하여 계산 및 제어 오버헤드를 줄이고, 큐 혼잡이나 링크 단절 위험, 라우팅 홀 등 위험 상황에서만 한정하여 기계학습 기반의 예측적 라우팅 분기(Predictive branch)를 selective하게 작동시킨다.

기존 \phaseTwelve\ 모델 대비 본 논문에서 제시하는 \method는 채널 유실 인지 스위칭, 상위-$k$ 전방향 안정성 지표, 에너지 고려 타이 브레이킹 및 패킷 드롭 억제 메커니즘을 추가하여 신뢰성과 물리 계층의 강인성을 대폭 강화하였다. 본 연구의 주요 기여는 다음과 같다.
\begin{itemize}
    \item 분산 FANET 라우팅 문제를 Dec-POMDP 모델로 공식화하고 오프라인 전역 학습 정보와 배포 시 로컬 관측 정보를 수학적으로 격리하여 로컬 실행 가능성을 보장하였다.
    \item 전역 교사의 행동 확률 분포의 결합적인 경향성을 1-hop 경량 학생 모델에 이전할 수 있도록 KL Divergence 및 Oracle 경로 안내 손실 함수를 포함하는 다중 목적 지식 증류 파이프라인을 제안하였다.
    \item 링크 잔여 수명, 버퍼 임계치, Onward 안정성을 종합한 위험도 점수(Danger Score)와 안전 게이트 메커니즘을 통해 위험 상황에서만 예측 모델을 구동하는 예측적 위험 스위칭 기법을 설계하였다.
    \item 총 14개 시나리오 및 5개 시드를 이용해 대규모 비교 실험을 수행하여, 제안 기법이 제어 오버헤드 바이트를 현저히 절약하면서도 GPSR, Predictive Geo, Evo-QGeo, IQMR, DRAMA 대비 PDR 및 데드라인 전송률 면에서 우수한 성과를 보임을 입증하였다.
\end{itemize}
